Item (9) also requires some explanation. In the annual accounts for 1850—1851, the fund's monetary assets on 1 July 1851 are entered as 125076 Thl. 1 Ggr. 11 Pf., of which the interest-bearing capital amounts to 123372 Thl. 2 Ggr. 7 Pf.; the remainder consists of cash on hand, arrears and an interest-free advance granted to the university pharmacist. About half the capital is invested in mortgage loans to private debtors, and the remainder in non-redeemable government securities. In the latest annual accounts, as in several earlier ones, these latter securities are entered simply at their face value. In 1845, by contrast, these entries were made at the purchase prices, and I retained them unchanged in my calculation at that time because fluctuations in the value of government securities were then much smaller than they have been over the last 3 to 4 years. Now that a considerable part of the government securities belonging to the widows' fund stands so far below face value, I consider it necessary to value the securities in the balance account at their current actual value, at which they can be realized. For this I used the stock-exchange prices in Frankfurt and Hanover on 1 October, since a definite date had to be chosen. They are as follows
\begin{center}
\footnotesize\setlength{\tabcolsep}{2pt}
\begin{tabular}{l r|l r}
Austrian \(4\tfrac{1}{2}\) Metalliques . .&\(67\tfrac{1}{8}\)&Hanoverian 5 per cent . . .&\(104\tfrac{1}{2}\)\\
— \(4\tfrac{1}{2}\) at Goll . . . .&\(69\tfrac{3}{4}\)&— 4 per cent . . .&102\\
— \(4\tfrac{1}{2}\) at Bethmann&\(71\tfrac{3}{4}\)&— \(3\tfrac{1}{2}\) Courant&\(98\tfrac{1}{4}\)\\
Baden \(3\tfrac{1}{2}\) . . . . . . . . . . . . . .&\(87\tfrac{3}{4}\)&— \(3\tfrac{1}{2}\) Gold . .&\(97\tfrac{3}{4}\)
\end{tabular}
\end{center}
\SourceLeaf{173}\noindent\hspace*{\parindent}For the gold rate I assumed 1 Louisd’or \(=\) 9 fl. 38 kr. and \(=5\tfrac{1}{2}\) thalers Courant. Some of the prices have, incidentally, fallen a little further since then, although only slightly.

The result of these conversions is that the same interest-bearing capital entered at a face value of 123372 thalers yields only 119044 Thl. 23 Ggr. at current stock-exchange prices. To this sum I add the 800 thalers for which the former Müller house was sold. Thus, with the other items added (cash on hand, etc.), the monetary assets of the widows' fund on 1 July 1851 come to
\begin{center}121548 thalers 22 Ggr.\end{center}
\noindent\hspace*{\parindent}To estimate the adjustment to 1 October 1851 as far as I am able, I subtract from this sum the widows' pensions that fell due at Michaelmas, amounting to . 1787 Thl. 12 Ggr., and add
\begin{center}
\small\setlength{\tabcolsep}{2pt}
\begin{tabular}{l r l r l}
\shortstack[l]{the contributions for the year 1850 to 1851\\. . . . . . . . . . . . . . . . . . . . . .}&780&,,&&\\
\shortstack[l]{the half-year's lease rent for the pharmacy\\. . . . . . . . . . . . . . . . . . . . . . . . . . .}&500&,,&&\\
\shortstack[l]{and a quarter-year's interest income,\\based on last year's income after deducting\\collection costs . . . . . . . . . . . . . . . . . . . . . .}&1248&,,&10&,,
\end{tabular}
\end{center}
\noindent which then gives item (9) above.

\begin{center}\rule{.28\linewidth}{.3pt}\end{center}
